%% kapa: exam=1 round=33 year=2022 period=2 form=A subject=accounting subject_ko=회계학 seq=30 number=70
%% source: Q-Net 감정평가사 기출문제 (cst00309) artlSeq=5209997 file=raw/1008/5209997/2022년도 제33회 감정평가사 1차시험 2교시(감정평가 관계법규, 회계학).pdf
%% conversion: kordoc 4.24.51 -> tools/gen1.py (auto); review=proofread-sonnet; flags=none
\begin{question}{70}
매각예정으로 분류된 비유동자산 또는 처분자산집단에 관한 설명으로 옳은 것은?
\begin{choices}
  \item 매각예정으로 분류하였으나 중단영업의 정의를 충족하지 않는 비유동자산(또는 처분자산집단)을 재측정하여 인식하는 평가손익은 계속영업손익에 포함한다.
  \item 소유주에 대한 분배예정으로 분류된 비유동자산(또는 처분자산집단)은 공정가치와 장부금액 중 작은 금액으로 측정한다.
  \item 비유동자산이 매각예정으로 분류되거나 매각예정으로 분류된 처분자산집단의 일부이더라도 그 자산은 감가상각 또는 상각을 중단하지 아니한다.
  \item 매각예정으로 분류된 비유동자산(또는 처분자산집단)은 공정가치와 장부금액 중 큰 금액으로 측정한다.
  \item 매각예정으로 분류된 처분자산집단의 부채와 관련된 이자와 기타 비용은 인식을 중단한다.
\end{choices}
\end{question}
\answer{1}
